\documentclass[10pt,a4paper]{amsart}
\usepackage[margin=19mm]{geometry}
\usepackage{amsmath,amssymb,mathtools}
\usepackage[scr=rsfs,cal=euler]{mathalfa}
\usepackage{xfrac}
\usepackage[unicode,hidelinks,bookmarks=false]{hyperref}
\newcommand{\ie}{i.e.\ }
\newcommand{\cf}{cf.\ }
\newcommand{\resp}{respectively\ }
\newcommand{\Subsection}{\subsection}
\providecommand{\nobreakdash}{\nobreak\hbox{-}}
\setlength{\emergencystretch}{2em}
\multlinegap0pt
\usepackage{amssymb}
\usepackage{mathtools}
\usepackage{tikz}
\usepackage{enumerate}
\usepackage{xcolor}
\newtheorem{proposition}{Proposition}
\newcommand{\sM}{{\sf M}}
\title{Second moments associated with finite free multiplicative convolution}
\author{Nicolas Gilliers, Andrea Pettenello}
\date{Source: arXiv:2608.24352v1 (CC BY 4.0)}
\begin{document}
\maketitle

\noindent\emph{Source and licence.} Second moments associated with finite free multiplicative convolution. Version arXiv:2608.24352v1 (CC BY 4.0), distributed by arXiv under the Creative Commons Attribution 4.0 International licence, http://creativecommons.org/licenses/by/4.0/, as recorded in the arXiv licence metadata for this exact version and shown on the abstract page https://arxiv.org/abs/2608.24352v1. Full licence notice: \texttt{licenses/arxiv-2608.24352-CC-BY-4.0.txt}.\par\smallskip

\noindent\textit{Editorial scope.} Counted unit: the article's proposition prop:sumovereta (source0.tex lines 572--582). The article's complete original proof of it is delivered with it (source0.tex lines 583--715, one \texttt{proof} environment of 133 lines). No mathematics is written, repaired or re-derived: the statement and the proof are the article's own lines, in the article's own notation, and no retained line is substituted or re-flowed.  No further source-local context is needed: every cross-reference of every retained line resolves inside the two delivered spans.  Nothing was omitted from any retained span, so the package carries no omission record.  No retained line contains a citation, so the wrapper carries no bibliography and no bibliography entry.  No retained line contains a figure environment, an image-inclusion command or drawing code, so no image is delivered and the package declares no asset dependency.  Editorial formatting only, in addition: an AMS \texttt{amsart} preamble that restates the article's own declarations and macros used by the retained lines, namely the environments \texttt{proposition} on separate counters, together with the standard packages those lines need.  The retained environments print in this document as Proposition 1 in order of appearance, whereas the article prints the same items with the numbering of its own full document.  The complete native source of the article as distributed in the arXiv e-print for this exact version is bundled unchanged as \texttt{sources/208465/source0.tex}.
\begin{proposition}\label{prop:sumovereta}
		Let $r\in [k]$. Fix $\gamma \leq {\sf K}$ a permutation less than ${\sf K} \in D(I,J)$. Under the condition $\iota(((I\cap J)_I)^{\gamma}) \,\Delta \,\iota(((I\cap J)_J)^{\gamma}) \neq \emptyset$, one has
	$$
		\sum_{\substack{\eta \leq {\sf M}\\ |\eta\gamma(I)\cap I|=r}}\varepsilon(\eta)=0.
	$$
	Otherwise, 
	\begin{align*}
		\sum_{\substack{\eta \leq {\sf M} \\ |\eta\gamma(I)\cap I|=r}}\varepsilon(\eta)=
		(-1)^{r-|I^\gamma|} \binom{|I\cap J|}{r-|(I\backslash I \cap J)^{\gamma}|}.
	\end{align*}
\end{proposition}

\begin{proof}
    We first define a partition of $(I \cap J)_I$ and a partition of $(I \cap J)_J$. Set
    \begin{align*}
		 & E^J \vcentcolon= \{ j \in (I\cap J)_J : (i,j) \in  \gamma \text{ for some } i \in I \},         \\
		 & L^J \vcentcolon= (I\cap J)_J^{\gamma},
	\end{align*}
	so that
    \begin{equation*}
        (I\cap J)_J = E^J\sqcup L^J,
    \end{equation*}
    and
     \begin{align*}
		 & E^I \vcentcolon= \{ i \in (I\cap J)_I : (i,j) \in  \gamma \text{ for some } j \in J \},         \\
		 & L^I \vcentcolon= (I\cap J)_I^{\gamma},
	\end{align*}
    so that 
    \begin{equation*}
        (I\cap J)_I = E^I \sqcup L^I.
    \end{equation*}
   Define ${\sf m} \in S_{I \sqcup J}$ by ${\sf m} = \prod_{x \in I \cap J}(x_I, x_J)$. Now consider the four block partition $P$ of $I \subset I \sqcup J$  given by
    \begin{equation*}
        P = \{E^I, L^I\} \wedge {\sf m}(\{E^J, L^J\}).
    \end{equation*}
    The blocks of the partition $P$ are indexed by pairs $(X,Y)$ where $X \in \{E^I,L^I\}$ and $Y \in \{E^J, L^J\}$ (the pair $(X,Y)$ corresponds to the block $X \cap {\sf m}(Y)$).
    Since $\eta$ ranges over the permutations less than $\sM$, its action is determined by its restriction to $(I\cap J)_I$. This restriction is, in turn, determined by the family $V^{\eta}_{(X,Y)} \subset (X,Y)$ of subsets on which it coincides with  ${\sf m}$ (since ${\sf m}$ is the largest permutation less than ${\sf M}$).
    Consider the partition of $I \subset I \sqcup J$ given by
    $$
    I =(I\backslash I\cap J)^{\gamma} \sqcup (I\backslash I\cap J)\backslash (I\backslash I\cap J)^{\gamma}\sqcup E^I \sqcup L^I.
    $$ We now partition the set $\mathcal{I}\vcentcolon= \eta\gamma(I) \cap I$ by intersecting it with the partition of $I$ defined above. Firstly, we have
	$$
		(I\backslash I\cap J)^{\gamma} \cap \mathcal{I}=(I\backslash I\cap J)^{\gamma},
	$$
	since $\eta$ acts trivially on elements in $I\backslash (I \cap J)$. Elements in $(I\backslash I\cap J)\backslash(I\backslash I\cap J)^{\gamma}$ are sent by $\gamma$ either to $(I\cap J)_J$ or to its complement in $J$. Elements falling in the latter set cannot be in $\mathcal{I}$ since $\eta$ acts trivially on their image by $\gamma$. Hence, elements $x \in \mathcal{I}\cap (I\backslash I\cap J)\backslash(I\backslash I\cap J)^{\gamma}$ fall in the former set and have images through $\gamma$ in the block $E^J$. That is, we have 
    $$
		\gamma(\mathcal{I}\cap(I\backslash I\cap J)\backslash(I\backslash I\cap J)^{\gamma}) \subset E^J.
	$$
    Applying $\eta$, we get that
	\begin{equation*}
		\eta\gamma(\mathcal{I}\cap(I\backslash I\cap J)\backslash(I\backslash I\cap J)^{\gamma}) \subseteq V^{\eta}(L^I,E^J)\sqcup V^{\eta}(E^I,E^J).
	\end{equation*}
    Elements in $E^I$ are mapped by $\gamma$ either to $E^J$ or to $J\setminus (I\cap J)_J$. Elements belonging to the latter set cannot be in $\mathcal{I}$ since $\eta$ acts trivially on their image by $\gamma$. Therefore, $\gamma$ necessarily maps elements in $\mathcal{I}\cap E^I$ to elements in the block $E^J$, which means that
    $$
		\eta\gamma(\mathcal{I}\cap E^I) \subseteq V^{\eta}(L^I,E^J)\sqcup V^{\eta}(E^I,E^J).
	$$
    By definition of the sets $(\cdot, E^J)$ and $V^{\eta}(\cdot,E^J)$, we have
    $$
    V^\eta(L^I,E^J)\subseteq \eta\gamma(\mathcal{I}\cap(I\backslash I\cap J)\backslash(I\backslash I\cap J)^{\gamma})\sqcup \eta\gamma(\mathcal{I}\cap E^I)
    $$
    and
    $$
    V^\eta(E^I,E^J) \subseteq \eta\gamma(\mathcal{I}\cap(I\backslash I\cap J)\backslash(I\backslash I\cap J)^{\gamma})\sqcup \eta\gamma(\mathcal{I}\cap E^I)
    $$
    Putting the above inclusions together, we find that
    $$
  \eta\gamma(\mathcal{I}\cap(I\backslash I\cap J)\backslash(I\backslash I\cap J)^{\gamma})\sqcup \eta\gamma(\mathcal{I}\cap E^I) =   V^\eta(L^I,E^J) \sqcup V^\eta(E^I,E^J).
    $$
    Finally, it is not difficult to see that
	\begin{equation*}
		\mathcal{I}\cap L^I= ((L^I,E^J)\backslash V^{\eta}(L^I,E^J))  \sqcup ((L^I,L^J)\backslash V^{\eta}(L^I,L^J)).
	\end{equation*}
	We therefore obtain 
	\begin{align*}
		\mathcal{I} = & (I\backslash I\cap J)^{\gamma} \sqcup ((L^I,E^J)\backslash V^{\eta}(L^I,E^J))       \sqcup ((L^I,L^J)\backslash V^{\eta}(L^I,L^J))                                                  \sqcup V^{\eta}(L^I,E^J)\sqcup  V^{\eta}(E^I,E^J).
	\end{align*}


	From the above partition we infer that
	\begin{align}
		\label{eqn:conditionone}
		 |V^{\eta}(E^I,E^J
        )|
		&= |\mathcal{I}| + |V^{\eta}(L^I,L^J)| \nonumber
		                                    -  (|(L^I,L^J)| + |(L^I,E^J)| + |(I\setminus (I\cap J))^\gamma|)                                        \nonumber \\
		                                   & = |\mathcal{I}| + |V^{\eta}(L^I,L^J)|-|I^\gamma|
	\end{align}
    Denote by $s$ the number of 2-cycles of $\eta$, so that
	\begin{align}
		\label{eqn:conditiontwo}
		\sum_{(X,Y)}|V^{\eta}(X,Y) | = s
	\end{align}
	From (\ref{eqn:conditionone}), it follows that the number of $\eta \leq {\sf M}$ with $s$ 2-cycles such that $|\mathcal{I}|=r$ is given by
	\begin{align}
	\label{eqn:CountingEta}
			 & \sum_v\binom{|(L^I,L^J)|}{v}
			 \binom{|(E^I, E^J)|}{r + v-|I^{\gamma}|} \binom{|(E^I,L^J)|+|(L^I,E^J)|}{s-r+|I^\gamma|-2v}
	\end{align}
    Observe that
    \begin{equation}
    \label{eqn:SymDiffEquality}
          |(E^I,L^J)| + |(L^I,E^J)| = |L^I|+ |L^J| -2|(L^I,L^J)|= |\iota(L^I)\Delta \iota (L^J)|
    \end{equation}
    
    Similarly, we have
    $$
    |(E^I,E^J)| = |E^I|-|(E^I,L^J)|= |E^I| - (|L^J|-|(L^I,L^J)|)  = |I\cap J| - |\iota(L^I) \cup \iota(L^J)|,
    $$
    where in the last equality we added and subtracted $|L^I|$.
    Finally, it is clear that
    $$
|(L^I,L^J)| = |\iota(L^I) \cap \iota (L^J)|.
    $$
    Using (\ref{eqn:CountingEta}) and the above equalities we conclude that
    \begin{align*}
		 & \sum_{\substack{\eta \leq {\sf M}                                                                                         \\ |\eta\gamma(I)\cap I|=r}}\varepsilon(\eta)= \sum_{s=0}^{|I\cap J|}\sum_v (-1)^{s-r-2v+|I^{\gamma}|}(-1)^{r-|I^{\gamma}|}
		  \binom{|\iota(L^I)\cap \iota(L^J)|}{v}\binom{|\iota(L^I) \Delta \iota(L^J)|}{s-r-2v+|I^{\gamma}|}  \\
		&\hspace{10cm}\times\binom{|I\cap J|-|\iota(L^I)\cup \iota(L^J)|}{r+ v-|I^{\gamma}|} ,
	\end{align*}
	Interchanging the sums over $s$ and $v$, we find
	\begin{align}
    \label{eqn:SumEta}
		 & \sum_{\substack{\eta \leq {\sf M}                                                                                        \nonumber\\ |\eta\gamma(I)\cap I|=r}}\varepsilon(\eta)= (-1)^{r-|I^{\gamma}|}\sum_v \binom{|\iota(L^I)\cap \iota(L^J)|}{v}\binom{|I\cap J|-|\iota(L^I)\cup \iota(L^J)|}{r+ v-|I^{\gamma}|}\\
		 & \hspace{5cm} \sum_{s=0}^{|I\cap J|} (-1)^{s-r-2v+|I^{\gamma}|}\binom{|\iota(L^I) \Delta \iota(L^J)|}{s-r-2v+|I^{\gamma}|}.
    \end{align}
     Under the condition $|\iota(L^I) \Delta \iota(L^J)|=0$, we get
	\begin{align*}
		\sum_{\substack{\eta \leq {\sf M} \\ |\eta\gamma(I)\cap I|=r}}\varepsilon(\eta)&= (-1)^{r-|I^{\gamma}|}\sum_v \binom{|\iota(L^I)\cap \iota(L^J)|}{v}\binom{|I\cap J|-|\iota(L^I)\cup \iota(L^J)|}{r+ v-|I^{\gamma}|}
	\end{align*}
    Suppose instead that $|\iota(L^I) \Delta \iota(L^J)|>0$. Using (\ref{eqn:conditionone}) and (\ref{eqn:SymDiffEquality}) we have 
    $$
    |I\cap J| -|V^\eta(E^I,E^J)|-|V^\eta(L^I,L^J)| \geq |I\cap J|- |(E^I,E^J)|- |(L^I,L^J)|= |\iota(L^I)\Delta \iota(L^J)|,
    $$
    which implies $|I\cap J|\geq |\iota(L^I) \Delta \iota(L^J)|+r+2v-|I^\gamma|$. Therefore, applying the binomial formula to the second sum in (\ref{eqn:SumEta}) yields
    $$
    \sum_{\substack{\eta \leq {\sf M} \\ |\eta\gamma(I)\cap I|=r}}\varepsilon(\eta) = 0.
    $$
    The desired result now follows from the Vandermonde identity upon noticing that
	$$
		|\iota(L^I)\cap \iota(L^J)|-|I^{\gamma}| =
		|\iota(L^I)|-|I^{\gamma}| = |L^I|-|I^{\gamma}| = |(I\backslash I\cap J)^{\gamma}|,
	$$
	under the condition that $\iota(L^I)\Delta \iota(L^J)=\emptyset$.
    
\end{proof}

\end{document}
